% method.tex — Section II: Method
% Target: PRD two-column (revtex4-2)
% Macros used: \HH, \dL, \Mz, \Mbh, \Mstellar, \pGW, \pdet, \pgal, \Lcat, \Lcomp, \SNR

We describe the end-to-end analysis pipeline: EMRI waveform generation,
LISA noise modeling, Fisher-matrix parameter estimation, and the dark siren
Bayesian inference framework used to constrain the Hubble constant.

%% ============================================================
\subsection{EMRI signal model}
\label{sec:emri_signal}
%% ============================================================

An extreme mass-ratio inspiral (EMRI) consists of a stellar-mass compact
object of mass $\mu \sim 1$--$100\,M_\odot$ spiraling into a massive black
hole (MBH) of mass $\Mbh \sim 10^4$--$10^7\,M_\odot$ under gravitational
radiation reaction.  The resulting millihertz gravitational-wave signal is
observable by the Laser Interferometer Space Antenna
(LISA)~\cite{LISA:2024hlh,Babak:2017tow} over observation times of order
years.

We generate EMRI waveforms using the augmented analytic kludge (AAK)
model~\cite{Chua:2020stf,Katz:2021yft}, which provides a fast, GPU-accelerated
template family capturing the dominant features of Kerr geodesic motion
including eccentricity, spin precession, and orbital inclination.
The LISA detector response is obtained by projecting the waveform onto
the time-delay interferometry (TDI) channels $A$ and $E$ via the
equal-arm-length approximation~\cite{Katz:2022yqe}.  The $T$ channel
is a null channel for the signal in this approximation and is neglected.

Each EMRI signal is parameterized by 14 source parameters:
the MBH mass $\Mbh$, compact-object mass $\mu$, dimensionless MBH spin $a$,
initial orbital semi-latus rectum $p_0$, initial eccentricity $e_0$,
inclination cosine $x_0 = \cos I$, luminosity distance $\dL$,
sky position angles $(\theta_S, \phi_S)$ in ecliptic coordinates,
spin orientation angles $(\theta_K, \phi_K)$,
and three initial orbital phases $(\Phi_{\phi 0}, \Phi_{\theta 0}, \Phi_{r0})$.
A summary of the parameter ranges is given in Appendix~\ref{app:parameters}.

The observation time is set to $T_\mathrm{obs} = 5\,\mathrm{yr}$ with a
sampling interval of $\Delta t = 10\,\mathrm{s}$, yielding
$\sim 1.6 \times 10^7$ time-domain samples per channel.

%% ============================================================
\subsection{LISA noise model and signal-to-noise ratio}
\label{sec:snr}
%% ============================================================

The noise power spectral density (PSD) for the $A$ and $E$ TDI channels
is~\cite{Babak:2023lro}
\begin{align}
  S_n^{(A,E)}(f) &= 8\sin^2\!\Bigl(\frac{2\pi f L}{c}\Bigr)
  \biggl[
    S_\mathrm{OMS}(f)\,\Bigl(\cos\frac{2\pi f L}{c} + 2\Bigr)
  \nonumber\\
  &\quad + 2\Bigl(3 + 2\cos\frac{2\pi f L}{c}
    + \cos\frac{4\pi f L}{c}\Bigr)\,S_\mathrm{TM}(f)
  \biggr]
  \nonumber\\
  &\quad + S_c(f) \,,
  \label{eq:psd}
\end{align}
where $L = 2.5\times10^9\,\mathrm{m}$ is the LISA arm length,
$c$ is the speed of light,
\begin{equation}
  S_\mathrm{OMS}(f) = (15\,\mathrm{pm})^2\,
  \Bigl(1 + \Bigl(\frac{2\,\mathrm{mHz}}{f}\Bigr)^{\!4}\Bigr)\,
  \Bigl(\frac{2\pi f}{c}\Bigr)^{\!2}
  \label{eq:s_oms}
\end{equation}
is the optical metrology noise, and
\begin{equation}
  S_\mathrm{TM}(f) = (3\,\mathrm{fm/s^2})^2\,
  \Bigl(1 + \Bigl(\frac{0.4\,\mathrm{mHz}}{f}\Bigr)^{\!2}\Bigr)
  \Bigl(1 + \Bigl(\frac{f}{8\,\mathrm{mHz}}\Bigr)^{\!4}\Bigr)
  \Bigl(\frac{1}{2\pi f\, c}\Bigr)^{\!2}
  \label{eq:s_tm}
\end{equation}
is the test-mass acceleration noise.
The term $S_c(f)$ is the residual galactic confusion noise from unresolved
white-dwarf binaries~\cite{Cornish:2017oic}, which depends on
the observation time and dominates the sensitivity in the
$0.1$--$3\,\mathrm{mHz}$ band.

The optimal signal-to-noise ratio (SNR) is defined as
\begin{equation}
  \SNR^2 = \sum_{\alpha \in \{A,E\}}
  4\,\mathrm{Re}\!\int_{f_\mathrm{min}}^{f_\mathrm{max}}
  \frac{|\tilde{h}^\alpha(f)|^2}{S_n^{\alpha}(f)}\,\mathrm{d}f \,,
  \label{eq:snr}
\end{equation}
where $\tilde{h}^\alpha(f)$ is the Fourier transform of the TDI response in
channel $\alpha$ and the integration extends from
$f_\mathrm{min} = 10^{-5}\,\mathrm{Hz}$ to
$f_\mathrm{max} = 1\,\mathrm{Hz}$.
We adopt a detection threshold of $\SNR \geq \SNR_\mathrm{th} = 20$.
As a computational optimization, a preliminary SNR estimate is performed
using $T_\mathrm{obs} = 1\,\mathrm{yr}$ of data; only events exceeding
$0.2\,\SNR_\mathrm{th}$ in this pre-screening proceed to the full
five-year waveform evaluation and Fisher-matrix computation.

%% ============================================================
\subsection{Parameter estimation}
\label{sec:fisher}
%% ============================================================

For detected events ($\SNR \geq \SNR_\mathrm{th}$), we estimate parameter
uncertainties via the Fisher information matrix~\cite{Vallisneri:2007ev}.
The Fisher matrix for parameters $\boldsymbol{\theta} = \{\theta^i\}$ is
\begin{equation}
  \Gamma_{ij} = \bigl\langle \partial_i h \,\big|\, \partial_j h \bigr\rangle \,,
  \label{eq:fisher}
\end{equation}
where the noise-weighted inner product $\langle \cdot | \cdot \rangle$
is defined by Eq.~\eqref{eq:snr} with the replacement
$|\tilde{h}^\alpha|^2 \to
\tilde{h}_1^{\alpha}(f)\,\tilde{h}_2^{\alpha *}(f)$,
and $\partial_i h \equiv \partial h / \partial \theta^i$.

Partial derivatives of the waveform are computed numerically using a
five-point central-difference stencil~\cite{Vallisneri:2007ev},
\begin{equation}
  \frac{\partial h}{\partial \theta^i} \approx
  \frac{-h(\theta^i + 2\epsilon) + 8\,h(\theta^i + \epsilon)
    - 8\,h(\theta^i - \epsilon) + h(\theta^i - 2\epsilon)}
  {12\,\epsilon} \,,
  \label{eq:five_point}
\end{equation}
with step sizes $\epsilon = 10^{-6}$ for each parameter (in the
parameter's natural units).
This $\mathcal{O}(\epsilon^4)$-accurate stencil requires four waveform
evaluations per parameter, for a total of $4 \times 14 = 56$ waveform
generations per Fisher matrix.  The resulting $14 \times 14$ Fisher matrix
$\Gamma_{ij}$ is symmetric by construction, so the inner product
computation exploits this symmetry.

The Cram\'er-Rao lower bound on the parameter covariance is obtained
by inverting the Fisher matrix,
\begin{equation}
  \Sigma_{ij} = \bigl(\Gamma^{-1}\bigr)_{ij} \,.
  \label{eq:cramer_rao}
\end{equation}
This provides a lower bound on the variance of any unbiased estimator and
is a reliable estimate of the true posterior width in the high-SNR
regime~\cite{Vallisneri:2007ev}.

For the dark siren analysis, the full $14 \times 14$ covariance matrix is
projected onto the subspaces of interest.
In the \emph{standard channel} (without MBH mass information), we retain the
$3 \times 3$ sub-block corresponding to $(\phi_S, \theta_S, \dL)$.
In the \emph{MBH mass channel}, we retain the $4 \times 4$ sub-block
$(\phi_S, \theta_S, \dL, \Mz)$, where $\Mz = \Mbh(1+z)$ is the
redshifted (detector-frame) MBH mass measured from the waveform.
These projected covariance matrices enter the gravitational-wave likelihood
described in Sec.~\ref{sec:likelihood}.

%% ============================================================
\subsection{Dark siren likelihood}
\label{sec:likelihood}
%% ============================================================

We follow the dark siren formalism of
Schutz~\cite{Schutz:1986gp}, extended by the galaxy catalog approach
of Refs.~\cite{Gray:2019ksv,Gair:2022zsa}, including the
completeness correction of Gray~et~al.~\cite{Gray:2019ksv}.

\subsubsection{Posterior on \texorpdfstring{$\HH$}{H0}}

Given $N_\mathrm{det}$ independent EMRI detections with data
$\{d_\mathrm{GW}^i\}$ and a galaxy catalog $\mathcal{C}$, the posterior
on the Hubble constant is
\begin{equation}
  p(\HH \mid \{d_\mathrm{GW}^i\}, \mathcal{C})
  \propto p(\HH)\,\prod_{i=1}^{N_\mathrm{det}}
  p(d_\mathrm{GW}^i \mid \HH, \mathcal{C}) \,,
  \label{eq:h0_posterior}
\end{equation}
where $p(\HH)$ is the prior on the Hubble constant, taken to be uniform
on $h \equiv \HH/(100\,\mathrm{km\,s^{-1}\,Mpc^{-1}}) \in [0.60, 0.86]$.

\subsubsection{Single-event likelihood}

The single-event likelihood marginalizes over the unknown host galaxy
identity by summing over all candidate host galaxies $j$ in the catalog.
We implement two channels that differ in the dimensionality of the
gravitational-wave measurement used.

\paragraph{Standard channel (without MBH mass).}
In this channel, the gravitational-wave data constrain only the sky
position and luminosity distance.  The single-event catalog likelihood for
detection $i$ is~\cite{Gray:2019ksv,Gair:2022zsa}
\begin{equation}
  \Lcat^{(i)} = \frac{%
    \sum_j \int \!\mathrm{d}z\;\pdet(z)\,
    \pGW\bigl(\phi_j, \theta_j, \dL(z, \HH)/\dL^\mathrm{det}\bigr)\,
    \pgal^{(j)}(z)
  }{%
    \sum_j \int \!\mathrm{d}z\;\pdet(z)\,\pgal^{(j)}(z)
  } \,,
  \label{eq:Lcat_no_mass}
\end{equation}
where $\pGW(\phi, \theta, \dL/\dL^\mathrm{det})$ is the three-dimensional
multivariate Gaussian likelihood constructed from the projected Fisher-matrix
covariance (Sec.~\ref{sec:fisher}) with mean at the measured values
$(\phi^\mathrm{det}, \theta^\mathrm{det}, 1)$;
$\pgal^{(j)}(z)$ is a Gaussian in redshift centered on the catalog value
$z_j$ with uncertainty $\sigma_{z,j}$;
$\dL(z, \HH)$ is the luminosity distance--redshift relation for the trial
value of $\HH$ in a flat $\Lambda$CDM cosmology~\cite{Hogg:1999ad};
and $\pdet(z)$ is the detection probability defined in
Sec.~\ref{sec:pdet}.
The numerator integrals are evaluated over the range
$\dL^\mathrm{det} \pm 4\sigma_{\dL}$ (converted to redshift at the trial
$\HH$), while the denominator integrals extend over
$z_j \pm 4\sigma_{z,j}$.
The sums run over all candidate host galaxies within the
$3\sigma$ localization volume of the detection.

\paragraph{MBH mass channel.}
The key idea of this work is to exploit the precise measurement of the
redshifted MBH mass $\Mz$ from the gravitational waveform as an additional
discriminant when identifying the host galaxy.  The detector-frame mass
is related to the source-frame mass by $\Mz = \Mbh(1+z)$.  Because the
Fisher-matrix uncertainty on $\Mz$ is extremely small
($\sigma_{\Mz}/\Mz \lesssim 10^{-6}$), the GW mass measurement acts
effectively as a delta function when convolved with the galaxy catalog mass
uncertainty ($\sigma_{\Mbh}/\Mbh \sim 20\%$).

To construct the MBH mass channel likelihood, we start from the
four-dimensional GW likelihood $\pGW(\phi, \theta, \dL/\dL^\mathrm{det},
\Mz/\Mz^\mathrm{det})$ built from the $4\times 4$ projected Fisher-matrix
covariance.  Marginalizing over the source-frame MBH mass $\Mbh$ with a
Gaussian catalog prior $g(\Mbh) = \mathcal{N}(\Mbh;\, \Mbh^{(j)},\,
\sigma_{\Mbh}^{(j)\,2})$ can be performed analytically.

% REVISED: presentation of analytic marginalization
Decomposing the four-dimensional GW covariance into a $3\times 3$ block
$\Sigma_\mathrm{obs}$ for $(\phi, \theta, \dL/\dL^\mathrm{det})$ and a
conditional distribution for $\Mz/\Mz^\mathrm{det}$ given the observed
variables, the marginalization over $\Mbh$ reduces to a one-dimensional
Gaussian integral.  Denoting the conditional mean and variance
of $\Mz/\Mz^\mathrm{det}$ by $\mu_\mathrm{cond}(\mathbf{x}_\mathrm{obs})$
and $\sigma^2_\mathrm{cond}$, respectively, and writing the catalog mass in
detector-frame fractional coordinates as
$\Mbh^{(j)}(1+z)/\Mz^\mathrm{det}$ with fractional uncertainty
$\sigma_{\Mbh}^{(j)}(1+z)/\Mz^\mathrm{det}$, the integral yields
\begin{align}
  &\int\!\mathrm{d}\Mbh \;\pGW(\mathbf{x}_\mathrm{obs}, \Mz/\Mz^\mathrm{det})
  \;g(\Mbh)
  \nonumber\\
  &\quad = \pGW^{(3)}(\mathbf{x}_\mathrm{obs})\,
  \frac{1}{\sqrt{2\pi(\sigma^2_\mathrm{cond} + \tilde{\sigma}_g^2)}}
  \nonumber\\
  &\qquad \times \exp\!\biggl(
    -\frac{(\mu_\mathrm{cond} - \tilde{\mu}_g)^2}
    {2(\sigma^2_\mathrm{cond} + \tilde{\sigma}_g^2)}
  \biggr) ,
  \label{eq:mass_integral}
\end{align}
where $\pGW^{(3)}$ is the three-dimensional marginal Gaussian,
$\tilde{\mu}_g = \Mbh^{(j)}(1+z)/\Mz^\mathrm{det}$, and
$\tilde{\sigma}_g = \sigma_{\Mbh}^{(j)}(1+z)/\Mz^\mathrm{det}$.
The single-event catalog likelihood for the MBH mass channel is then
\begin{equation}
  {L_\mathrm{cat,\Mz}}^{(i)} = \frac{%
    \sum_j \int\!\mathrm{d}z\;
    \pdet(z,\Mbh)\,
    \pGW^{(3)}\, \mathcal{G}_j(z)\,
    \pgal^{(j)}(z)
  }{%
    \sum_j \int\!\mathrm{d}z\!\int\!\mathrm{d}\Mbh\;
    \pdet(z,\Mbh)\,\pgal^{(j)}(z)\,g(\Mbh)
  } \,,
  \label{eq:Lcat_mass}
\end{equation}
where $\mathcal{G}_j(z)$ denotes the Gaussian factor from
Eq.~\eqref{eq:mass_integral} and the denominator accounts for the
two-dimensional selection effect in $(z, \Mbh)$.

\subsubsection{Completeness correction}

The galaxy catalog is incomplete, particularly at higher redshifts.
Following Gray~et~al.~\cite{Gray:2019ksv}, the full single-event likelihood
$p(d_\mathrm{GW}^i \mid \HH, \mathcal{C})$ is a weighted combination of the
catalog term $\Lcat$ and a completion term $\Lcomp$ that accounts for
galaxies absent from the catalog:
\begin{equation}
  p(d_\mathrm{GW}^i \mid \HH, \mathcal{C})
  = f_i\,\Lcat^{(i)} + (1 - f_i)\,\Lcomp^{(i)} \,,
  \label{eq:completeness_combination}
\end{equation}
where $f_i \in [0,1]$ is the catalog completeness fraction evaluated at the
inferred redshift of detection $i$ for the trial value of $\HH$.
The completion term assumes a uniform-in-comoving-volume galaxy
distribution,
\begin{equation}
  \Lcomp^{(i)} = \frac{%
    \int \!\mathrm{d}z\;\pGW(z)\,\pdet(z)\,
    \frac{\mathrm{d}V_c}{\mathrm{d}z}
  }{%
    \int \!\mathrm{d}z\;\pdet(z)\,
    \frac{\mathrm{d}V_c}{\mathrm{d}z}
  } \,,
  \label{eq:Lcomp}
\end{equation}
where $\mathrm{d}V_c/\mathrm{d}z$ is the comoving volume element
and $\pGW(z)$ denotes the three-dimensional GW likelihood marginalized
over sky position, evaluated along the line of sight of the detection.
For the MBH mass channel, the completion term uses the same
three-dimensional likelihood as Eq.~\eqref{eq:Lcomp}, because an
uncataloged host galaxy has no associated mass estimate.

%% ============================================================
\subsection{Galaxy catalog}
\label{sec:catalog}
%% ============================================================

We use the GLADE+ galaxy catalog~\cite{Dalya:2021ewn}, which provides
nearly full-sky coverage with approximately 23 million objects.
For each galaxy, we extract the redshift $z$, redshift uncertainty
$\sigma_z$, stellar mass $\Mstellar$, stellar mass uncertainty
$\sigma_{\Mstellar}$, and sky position $(\theta, \phi)$.
We restrict the catalog to galaxies with spectroscopic or
photometric redshifts and require that $\Mstellar$ is available.

The central MBH mass is estimated from the stellar mass via the
scaling relation of Reines~\&~Volonteri~\cite{Reines:2015moa},
\begin{equation}
  \log_{10}\!\Bigl(\frac{\Mbh}{M_\odot}\Bigr)
  = \alpha + \beta\,\log_{10}\!\Bigl(\frac{\Mstellar}{10^{11}\,M_\odot}\Bigr) ,
  \label{eq:scaling_relation}
\end{equation}
with $\alpha = 7.45 \pm 0.08$ and $\beta = 1.05 \pm 0.11$.
The uncertainty on the inferred $\Mbh$ is propagated from
$\sigma_{\Mstellar}$ and the intrinsic scatter of the scaling relation,
yielding typical fractional uncertainties
$\sigma_{\Mbh}/\Mbh \sim 20\%$.

Candidate host galaxies for each detection are identified by a ball-tree
spatial query within the $3\sigma$ error volume defined by the Fisher-matrix
uncertainties on $(\phi_S, \theta_S, \dL)$.
We restrict the MBH mass range to
$\Mbh \in [10^5,\, 10^6]\,M_\odot$
to match the mass range covered by the injection campaign
(Sec.~\ref{sec:pdet}).

%% ============================================================
\subsection{Detection probability}
\label{sec:pdet}
%% ============================================================

The detection probability $\pdet(z, \Mbh, \HH)$ is required both in the
likelihood numerator and in the selection-effect correction
(denominators of Eqs.~\eqref{eq:Lcat_no_mass}--\eqref{eq:Lcomp}).
We estimate it via a Monte Carlo injection campaign.

A two-dimensional grid in $(h, z)$ is constructed with 15 values of
$h \in [0.60, 0.86]$ and 10 redshift bins from $z = 0$ to $z_\mathrm{max}$.
For each grid cell, EMRI events are injected with fixed $(h, z)$ and
randomized extrinsic parameters (sky position, spin orientation,
orbital phases, compact-object mass, eccentricity, inclination).
For each injection, the full waveform is generated, the LISA TDI response
is computed, and the SNR is evaluated; a detection is counted when
$\SNR \geq \SNR_\mathrm{th}$.
The detection fraction in each bin provides the empirical estimate of
$\pdet$.

For continuous evaluation at arbitrary $(h, z)$, the grid is interpolated
using a regular-grid interpolator.  Per-bin uncertainties are quantified
using Wilson confidence intervals~\cite{Wilson:1927}.
At grid cells near the detection boundary where the detection fraction
transitions from zero to one, variance reduction is achieved through
importance-sampling-enhanced
estimation~\cite{Tiwari:2017ndi} with Neyman allocation, yielding
variance reduction factors of $11.8$--$24.9\times$ in the boundary
bins~\cite{Farr:2019rap}.

%% ============================================================
\subsection{Simulation procedure}
\label{sec:procedure}
%% ============================================================

We adopt a flat $\Lambda$CDM cosmology with fiducial parameters
$h_\mathrm{true} = 0.73$ and $\Omega_m = 0.25$.
The analysis proceeds as follows for each trial value of
$h \in [0.60,\, 0.86]$ with spacing $\Delta h \leq 0.01$:

\begin{enumerate}
  \item A galaxy is drawn from the GLADE+ catalog restricted to the MBH
    mass range $[10^5,\, 10^6]\,M_\odot$.

  \item The luminosity distance is computed from the galaxy redshift
    using the $\dL$--$z$ relation at $h_\mathrm{true}$.

  \item An EMRI waveform is generated with the galaxy's sky position,
    MBH mass, and luminosity distance, supplemented by randomized
    extrinsic parameters.

  \item If $\SNR \geq \SNR_\mathrm{th}$, the Fisher information matrix
    is computed and inverted to obtain the Cram\'er-Rao covariance matrix.
    Only events with $\sigma_{\dL}/\dL < 10\%$ are retained.

  \item For each retained detection, the $\HH$ posterior
    [Eq.~\eqref{eq:h0_posterior}] is evaluated by summing the
    single-event likelihood over all candidate host galaxies identified
    from the catalog.  Both the standard and MBH mass channels are
    evaluated for each detection.  The completeness-corrected likelihood
    [Eq.~\eqref{eq:completeness_combination}] is used throughout.
\end{enumerate}

The posterior is evaluated on a grid in $h$ and the combined constraint
from all detected events is obtained by multiplying the per-event
likelihoods [Eq.~\eqref{eq:h0_posterior}].

%% CITATIONS NEEDED (for gpd-bibliographer):
%% - MISSING:cornish-robson-2017 — Cornish & Robson (2017), arXiv:1703.09858, galactic confusion noise model
%% - MISSING:wilson-1927-intervals — Wilson (1927), confidence interval for binomial proportion
